\section*{Podsumowanie}

Przegląd przedstawionych architektur systemów RAG ujawnia wyraźną progresję od prostych, statycznych implementacji do wysoce adaptacyjnych i inteligentnych systemów zdolnych do autonomicznego podejmowania decyzji. 
Każda kolejna generacja rozwiązuje fundamentalne ograniczenia poprzedników, wprowadzając jednocześnie nowe możliwości i wyzwania techniczne.

Ewolucja od Naive RAG do Advanced RAG koncentrowała się na poprawie jakości retrieval poprzez optymalizacje przed i po etapie wyszukiwania, zachowując przy tym sekwencyjną naturę przetwarzania. 
Przejście do Modular RAG przełamało te ograniczenia, wprowadzając elastyczność architekturalną i możliwość adaptacji do różnorodnych scenariuszy zastosowań. 
Najnowsze systemy Agentic RAG reprezentują kolejny skok jakościowy, wprowadzając elementy autonomii i samooceny, które przybliżają systemy AI do bardziej naturalnych i inteligentnych form interakcji z wiedzą.

Ta ewolucja odzwierciedla szerszy trend w dziedzinie sztucznej inteligencji, gdzie statyczne systemy ustępują miejsca adaptacyjnym agentom zdolnym do uczenia się, planowania i autonomicznego podejmowania decyzji. 
Systemy RAG nowej generacji nie tylko lepiej wykorzystują zewnętrzne źródła wiedzy, ale także rozwijają się w kierunku bardziej złożonych i inteligentnych interakcji, które mogą znacząco poprawić jakość i efektywność rozwiązywania problemów w różnych dziedzinach.